\documentclass[12pt,dvipsnames]{amsart}
\usepackage{loom}
%\usepackage[T1]{fontenc}
\usepackage{microtype}        % better justification, fewer overful hboxes

%%%%%%%%%%%%%%%%%%
%% Font Choices
%%%%%%%%%%%%%%%%%%

%\usepackage{lmodern}          % safe default, good screen rendering
%\usepackage{mathpazo}         % Palatino: wider, more readable than Times
%\usepackage{fourier}          % Utopia: excellent for math-heavy documents

% line spacing in the document
\linespread{1.2}

% xcolor commonly causes option clashes, this fixes that
\PassOptionsToPackage{dvipsnames,table}{xcolor}
\usepackage[tmargin=1in, bmargin=1in, lmargin=1.3in, rmargin=1.3in]{geometry}

% included exclusively for clickable table of contents
\usepackage{hyperref}
\hypersetup{
    colorlinks,
    citecolor=black,
    filecolor=black,
    linkcolor=black,
    urlcolor=black
}

% use the better epsilon
\renewcommand{\epsilon}{\varepsilon}

% suppress the warning page
\pdfsuppresswarningpagegroup=1

% enable synctex for inverse search
\synctex=1

% include packages
\usepackage{base-macros,math-env,quiver}
\usepackage{microtype}

% colored text shortcuts

\newcommand{\isaac}[1]{\textcolor{Mahogany}{(Isaac) #1}}
\newcommand{\blue}[1]{\color{MidnightBlue}{#1}}
\newcommand{\red}[1]{\textcolor{Mahogany}{#1}}
\newcommand{\green}[1]{\textcolor{ForestGreen}{#1}}

% use mathptmx pkg while using default mathcal font
\DeclareMathAlphabet{\mathcal}{OMS}{cmsy}{m}{n}

% fixes the positioning of subscripts in $$ $$
\renewcommand{\det}{\operatorname{det}}

% number the pages
\pagenumbering{arabic}

\usetikzlibrary{positioning, arrows.meta}
\newcommand{\here}[2]{\tikz[remember picture]{\node[inner sep=0](#2){#1}}}



% bibliography
%\usepackage[citestyle=alphabetic,bibstyle=authortitle]{biblatex}
\usepackage[backend=biber,style=alphabetic,sorting=anyt]{biblatex}
\usepackage[new]{old-arrows} % for \longhookrightarrow
\addbibresource{refs.bib}

% document specific commands
\newcommand{\vir}{\text{vir}}

% checklist
\newlist{todolist}{itemize}{2}
\setlist[todolist]{label=$\square$}

\title{Relative Virtual Localization}
\begin{document}
\maketitle
% \pagestyle{empty}
\bigskip

\begin{center}
\today
\end{center}

\tableofcontents

\newpage
A proposal for a restructured document outline:
\begin{itemize}
  \item \textbf{Introduction.} Make the rest of the paper readable by someone who knows GP99 and Manolache but not Romagny fixed stacks. State the theorem, provide historical context, name the failure of torus localization on Artin bases.
    \begin{itemize}
      \item \textbf{Section 1.1} The formula, main theorem and what is new
      \item \textbf{Section 1.2} Three different failures of the Artin base: \(\iota_*\) is not invertible (may not even exist), \(\mathcal Y^T\) is not quasi-compact or finite type necessarily (\(B\mathbb G_m\)), \(\mathcal Y^T\) not algebraic (Tate curve).
      \item \textbf{Section 1.3} Reader's guide
            \textbf{Section 1.4} Conventions (cohomological vs. homological grading, Euler class conventions), standing hypotheses, notation.
    \end{itemize}
  \item \textbf{Section 2.} The fixed stack and its deformation theory
        \begin{itemize}
          \item \textbf{Section 2.1} Group actions on Artin stacks, the different fixed stacks. Prove relationships in the appendix but state relationships here. Good opportunity to highlight reparameterization.
          \item \textbf{Section 2.2} Properties of \(\eta\)
          \item \textbf{Section 2.3} \(\tau_{\geq -1}L_\eta\) is moving + the fixed relative obstruction theory
          \item \textbf{Section 2.4} Failure of algebraicity, the Tate curve example, explanation of why it is/isn't a problem
        \end{itemize}
  \item \textbf{Section 3.} The residue map. EG motivation for the construction written out in the first part of the section.
        \begin{itemize}
          \item \textbf{Section 3.1} The localized equivariant Chow groups and Euler classes of vector bundle stacks
          \item \textbf{Section 3.2} The construction of the residue, existence, independence, definition, cone-stack counterexample
          \item \textbf{Section 3.3} Compatability with Edidin-Graham: \textbf{Res} is the inverse when one exists. Just a sanity check.
        \end{itemize}
  \item \textbf{Section 4.} Localization formula
        \begin{itemize}
          \item \textbf{Section 4.1} The localization formula
          \item \textbf{Section 4.2} Specialization to GP99
        \end{itemize}
  \item \textbf{Section 5.} Application to stable log maps, as much as makes sense.
        \begin{itemize}
          \item \textbf{Section 5.1} Can at least verify the hypotheses in the setup, explain in broad strokes the setup
          \item \textbf{Section 5.2} Calculations where possible
        \end{itemize}
        \item \textbf{Appendix A.} Comparisons of fixed stacks, especially the na\"ive fixed locus. Proofs live here, statements originally encountered in Section 2.
\end{itemize}

\newpage

Hypotheses:
\begin{itemize}
  \item \(\mathcal Y\) or \(\mathcal X\) is an algebraic stack of finite type over a field \(k\) with the action of a torus \(T\), \(\eta:\mathcal Y^T\to \mathcal Y\) the canonical map, \(\mathcal Z\subseteq \mathcal Y^T\) is open and closed substack
  \item \(T\) is a split torus, or \(k = \overline k\). This is needed to ensure \(e_T(V)\) is a unit whenever \(V\) is finite-rank moving
  \item \(\mathcal Y\) is stratified by global quotients, needed for Kresch's Chow groups and Manolache's virtual pullback.
  \item \(\mathcal Z\) has the resolution property, needed for the existence of a POT for \(\eta:\mathcal Z\to \mathcal Y\).
  \item \(\mathcal Z\) is algebraic \(\iff\) \(\eta\) representable by algebraic spaces. Also want \(\mathcal Z\) to be finite type over \(k\) and to contain the image of \(\pi^T:\mathcal M^T\to \mathcal Y^T\).
\end{itemize}

Known reivions that must be made:
\begin{todolist}
  \item Prove main technical lemma, that \((\tau_{\geq -1}L_\eta)^f = 0\).
  \item Choose either \(\mathcal X\) or \(\mathcal Y\) as consistent notation for the base
  \item State the main localization theorem in a complete way
  \item Expand the section on the non-algebraicity of the fixed stack to include more exposition, currently only an example
  \item Find homes for everything in the ``Remaining facts from Tom's note'' section
  \item Ensure the proof of the \(e_T(N^{vir})\cap \Res_{M^T/\mathcal Z}(\widetilde \pi^!_{\widetilde E}(\alpha)) = (\pi^T)^f_{E}(\alpha)\) is complete enough
  \item Write Introduction
  \item Reorganize section 2
  \item Reorganize section 3
  \item Reorganize section 4
  \item Write section 5/decide what to include there
  \item Write new subsection in 3 verifying that $\Res$ recovers the inverse of \(\iota_*\) for EG localization
  \item Write new subsection in 4 verifying that after reparameterization this localization theorem recovers the old GP formula
\end{todolist}

\newpage

\section{Introduction}\label{rl-0001}

I don't want to write an introduction right now, but here's some stuff to mention.
\begin{itemize}
  \item We use the Chow groups of Kresch, but \(\hat A_*^T\mathcal Y\) does not mean the Chow groups of Edidin-Graham-Totaro that appear in his Definition 2.1.4 (iii). We need this group in order to make sense of the Euler classes of vector bundle stacks.
  \item The perfect obstruction theory of \(M^T\) is relative to the map \(M^T\to \mathcal Y^T\), \textit{not} the map \(M^T\to \mathcal Y\). We get this obstruction theory by restricting \(E\) to \(M^T\) and taking the fixed portion. This is the primary source of difficulty. We must find a way to relate the virtual fundamental classes \([M]^{vir}\) and \([M^T]^{vir}\) by studying the map \(\eta:\mathcal Y^T\to \mathcal Y\) between the bases, which is impossible via traditional localization theory since the main localization theorem fails for Artin stacks. We fix this by defining a residue map from classes of \(\mathcal Y\) to classes of \(\mathcal Y^T\). There is no canonical obstruction theory for \(\eta\) so the virtual pullback \(\eta^!\) (if it exists) depends on the choice of POT chosen, but this dependency can be fixed by capping with an appropriate inverse Euler class.
\end{itemize}

Each section begins with its own introduction explaining its context within the broader goal of the paper and is followed by subsections containing the actual math.

\bigskip

Throughout \(\mathcal X\) is an algebraic stack of finite type over a field \(k\) with an action of \(T\) a torus. We'll always denote by \(\eta:\mathcal X^T\to \mathcal X\) the canonical map from the Romagny fixed locus to \(\mathcal X\), and \(\mathcal Z\subseteq \mathcal X^T\) will be an open and closed sublocus. The restriction of \(\eta\) to \(\mathcal Z\) will also just be denoted \(\eta\). We'll need the following hypotheses.

\noindent
\begin{itemize}
  \item \(T\) is a split torus. We need this so that \(e_T(V)\) is a unit whenever \(V\) is a finite-rank moving \(T\)-equivariant locally free sheaf on \(\mathcal Z\subseteq \mathcal X^T\).
  \item \(\mathcal X\) is stratified by global quotients. Need this for Kresch's Chow groups and also for Manolache's virtual pullback.
  \item \(\mathcal Z\) has the resolution property. We need this for the existence of a POT for \(\eta:\mathcal Z\to \mathcal X\). Note this implies \(\mathcal Z = [U/\GL_n]\) with \(U\) some quasi-affine by \cite{totaro_ResolutionPropertySchemes2004}.
  \item \(\mathcal X^T\) is algebraic. This is not automatic, and is equivalent to \(\eta:\mathcal X^T\to \mathcal X\) being representable by algebraic spaces.
\end{itemize}

This means \(A^T_*(\mathcal Z) = A_*(\mathcal Z\times BT)\) exists with no further hypotheses.

\section{The fixed stack and its deformation theory}\label{rl-0002}

The aim of this section is to recall the definition of the Romagny fixed stack and establish some preliminary results regarding its deformation theory. Throughout this section \(G\) is a group scheme acting on an algebraic stack \(\mathcal X\).

\subsection{Preliminaries}\label{rl-0003}
Group actions on algebraic stacks are subtle, and there are several distinct notions of ``the fixed stack'' for a group acting on a stack. We make no attempt at a comprhensive treatment here; instead, we review 

We begin by reviewing group actions on stacks as defined by Romagny \cite{romagny_GroupActionsStacks2005} and four distinct notions of ``the fixed stack of a group action on a stack''. Throughout the main body of this paper we use the Romagny fixed stack

\begin{definition}\label{rl-0004}
\cite[Definition 2.1]{romagny_GroupActionsStacks2005}
  \label{defn:group-action-on-a-stack}
  Let \(\mathcal X\) be a stack over a scheme \(S\) and let \(G\) be a group algebraic space. Let \(m:G\times G\to G\) be the multiplication of \(G\) and \(e:S\to G\) its unit section. Let \(U\) be an \(S\)-scheme.
  \begin{enumerate}[(i)]
    \item An \textit{action} of \(G\) on \(\mathcal X\) is a morphism of stacks \(\mu:G\times \mathcal X\to \mathcal X\) satisfying the following 1-commutative diagrams
          \begin{equation}\label{eqn:group-action-diagrams}
            \begin{tikzcd}
            G\times G\times \mathcal X \arrow[r, "m\times \id_{\mathcal X}"] \arrow[d, swap, "\id_G\times \mu"] & G\times\mathcal X \arrow[d, "\mu"] \\
            G\times \mathcal X \arrow[r, "\mu"] & \mathcal X
            \end{tikzcd}\hphantom{asdfasdf}
            \begin{tikzcd}
            G\times \mathcal X \arrow[r,"\mu"] & \mathcal X \\
            \mathcal X\arrow[u, "e\times \id_{\mathcal X}"] \arrow[ur, swap, "\id_{\mathcal X}"{name=D}] &
            \end{tikzcd}
          \end{equation}
          The pair \((\mathcal X, \mu)\) is called a \(G\)-stack.
    \item A \(1\)-morphism of \(G\)-stacks or a \emph{\(G\)-equivariant morphism} between \((\mathcal X, \mu)\) and \((\mathcal Y, \nu)\) is a pair \((f, \sigma)\) where \(f:\mathcal X\to \mathcal Y\) is a \(1\)-morphism and \(\sigma:\nu\circ (\id_G\times f)\to f\circ \mu\) is a \(2\)-isomorphism written on sections as \(\sigma^x_g:g\cdot f(x) \xrightarrow{\sim} f(g\cdot x)\), such that \(\sigma_g^{h\cdot x}\circ g \cdot \sigma^x_h = \sigma^x_{gh}\) for all \(U\), \(x\in \mathcal X(U)\) and \(g, h\in G(U)\).
    \item A \textit{2-morphism} of \(G\)-stacks between 1-morphisms \((f_1,\sigma_1)\) and \((f_2, \sigma_2)\) is a 2-morphism \(\tau:f_1\Rightarrow f_2\) so that for all \(x\in \mathcal X(U)\) and \(g\in G(U)\) we have \(\sigma_{2,(g,x)}\circ g\cdot \tau_x = \tau_{g\cdot x}\circ \sigma_{1,(g,x)}\); by \(g\cdot \tau_x\) here we mean \(\nu(g, \tau_x)\).
    \item An \textit{isomorphism} of \(G\)-stacks is a \(1\)-G-morphism that is also an equivalence of groupoids over \(S\). A monomorphism is a fully faithful \(1\)-G-morphism and an epimorphism is a \(1\)-G-morphism that is locally essentially surjective.
  \end{enumerate}
\end{definition}

\begin{remark}\label{rl-0005}
  Technically what is stated here is what Romagny calls a \emph{strict action}; the more general version asks that these diagrams only commute up to 2-isomorphism. However Romagny immediately shows that strictification is a 2-equivalence \cite[Proposition 1.5]{romagny_GroupActionsStacks2005}
\end{remark}


Consider for a moment the fixed locus \(X^G\) of a scheme \(X\) with a group action \(G\). A map \(U\to X^G \in X^G(U)\) from some scheme \(U\) is exactly the same data as a \(G\)-equivariant map \(U\to X\) where \(U\) is endowed with the trivial action of \(G\). Romagny defines the fixed stack analogously to this functor of points description of \(X^G\). 
\begin{definition}\label{rl-0006}
  \label{defn:romagny-fixed-stack}
  Let \(G\) be a group algebraic space over \(S\) a scheme and let \(\mathcal X\) be a \(G\)-stack over \(S\). For an \(S\)-stack \(Y\), denote by \(\iota(\mathcal Y)\) the \(G\)-stack over \(S\) obtained by endowing \(\mathcal Y\) with the trivial \(G\)-action. The \textit{stack of fixed points} \(\mathcal X^{G}\) is a stack that \(2\)-represents the functor
  \[\sHom_{G-\St}(\iota(-), \mathcal X).\]
  Note that \(\mathcal X^G\) is an object in the category of stacks over \(S\), and \textit{not} in the category of \(G\)-stacks over \(S\). There is a canonical morphism \(\eta:\mathcal X^G\to \mathcal X\) which sends \((f,\sigma)\in \mathcal X^G(U)\) to \(f\in \mathcal X(U)\), where here we consider \(\mathcal X\) without its \(G\)-stack structure. 
\end{definition}

\begin{remark}\label{rl-0007}\label{rmk:description-of-homotopy-fixed-stack}
  Here is a slightly more explicit description of the points of $\mathcal X^{G}$ over a scheme $U \in \Sch/S$. An object of $\mathcal X^{G}(U)$ consists of a $G$-equivariant morphism $(x,\alpha):\iota(U)\to X$. The equivariance diagram in this case is
  \begin{center}
    \begin{tikzcd}[column sep={1.5cm,between origins},
      row sep={1.5cm,between origins}]
      G\times U\arrow[r, "\pr_2"] \arrow[d, swap, "\id_G \times x"] & U \arrow[d, "x"] \\
      G\times \mathcal X\arrow[r, "\mu"] \arrow[ur, Rightarrow, "\alpha"]& \mathcal X
    \end{tikzcd}
  \end{center}
  since $\iota(U)$ is simply $U$ with the trivial $G$-action. The morphism $x$ is an object of \(\mathcal X\) by the Yoneda lemma and the 2-commutativity implies that for all $g \in G(U)$ we get an isomorphism $\alpha_g:g\cdot x\xrightarrow{\sim}x$. Similarly, given a pair $(x,\{\alpha_g\}_{g\in G(U)})$ consisting of an object $x \in \mathcal X(U)$ and a collection of isomorphisms $\alpha_g:gx\xrightarrow{\sim}x$, we obtain a $G$-equivariant map $\iota(U)\to \mathcal X$. Thus the objects of $\mathcal X^{G}(U)$ are precisely pairs $(x,\{\alpha_g\}_{g\in G(U)})$ where $x \in \mathcal X(U)$ and the $\alpha_g:g\cdot x\xrightarrow{\sim} x$ are isomorphisms satisfying the cocycle condition \(\alpha_{gh} = \alpha_g\circ g\cdot \alpha_h\) pictured by
  \begin{equation}\label{eqn:cocycle-condition-for-fixed-stack}
    \begin{tikzcd}[column sep={1.5cm,between origins},
      row sep={1.5cm,between origins}]
      & gx \arrow[dr,"\alpha_g"] & \\
      (gh)x\arrow[ur, "g\cdot \alpha_h"] \arrow[rr, "\alpha_{gh}"] & & x
    \end{tikzcd}
  \end{equation}
  for all $g, h\in G(U)$. The morphism \(\eta:\mathcal X^G\to \mathcal X\) simply forgets the data of \(\alpha\).

  In particular, when \(G\) acts trivially on \(\mathcal X\), \(gx = x\) for all \(x\in \mathcal X(U)\) and the cocycle condition reduces to the requirement that \(\alpha_{gh} = \alpha_g\circ \alpha_h\); i.e. that \(\alpha:G\to \underline{\Aut}(x)\) is a group homomorphism. Thus, for the trivial action, the fiber of \(x\in \mathcal X(U)\) over \(\eta\) is the sheaf \(\underline{\Hom}_{S-gp}(G, \Aut(x))\).
\end{remark}


When there are many choices of \(\alpha_g\) for a single \(x\in \mathcal X(U)\) we get multiple points in \(\mathcal X^G(U)\) mapping to a single point in \(\mathcal X\). This means that \(\eta:\mathcal X^G\to \mathcal X\) is, in general, not a monomorphism.

\begin{example}\label{rl-0008}
  Let \(\mathfrak M\) denote the stack of nodal curves of some fixed genus (we suppress little \(g\) from the notation since we want to use it to denote a group element) and endow it with the trivial action of an algebraic group \(G\). Consider a scheme \(B\). Specializing Remark \ref{rmk:description-of-homotopy-fixed-stack} to this situation, we see that the points of \(\mathfrak M^G(B)\) are given by pairs \((x,\alpha)\) where \(x\in \mathfrak M(B)\) and \(\alpha:G(B)\to \Aut(x)\) is a group homomorphism. The point \(x\) defines a family of nodal curves \(E\to B\) over \(B\) by pulling back the universal curve over \(\mathfrak M\), and \(\alpha\) gives an automorphism
  \begin{center}
    \begin{tikzcd}
      E\arrow[r, "\sim"] \arrow[rd] & E \arrow[d]\\
      & B
    \end{tikzcd}
  \end{center}
  of this family for every element of \(G(B)\). Thus, \(\mathfrak M^G\) parameterizes families of nodal curves with an action of \(G_B\) on \(E\) over \(B\), and \(\eta:\mathfrak M^G\to \mathfrak M\) is the map that forgets this \(G\)-action. Whenever there is such a family \(E\to B\) with a nontrivial \(G_B\)-action on \(E\), \(\eta\) will not be a monomorphism.
\end{example}


\begin{example}\label{rl-0009}\label{example:BGm-and-its-fixed-stack}
  Let \(S\) be a scheme with a trivial \(\mathbb G_m\) action, and endow \(\mathcal X = [S/\mathbb G_m] = B\mathbb G_m\) with the trivial \(T = (\mathbb G_m)^n\) action. Every point in \(x \in \mathcal X\) is fixed by \(T\), so a collection \(\{\alpha_t\}\) satisfying the cocycle condition \eqref{eqn:cocycle-condition-for-fixed-stack} is the same data as a homomorphism \(\alpha:T\to \Aut_{\mathcal X}(x) = \mathbb G_m\), i.e. a character of the torus. The fixed stack \(\mathcal X^T\) is therefore a countable union of copies of \(\mathcal X\):
  \[\mathcal X^T = \coprod_{\chi\in \mathbb Z^n} \mathcal X,\]
  and the fiber of \(\eta\) at a point \(x \in \mathcal X\) is \(\mathbb Z^n\). In particular, \(\eta\) is not a monomorphism and is not quasi-compact. The pushforward \(\eta_*\), the thing which would becomes an isomorphism in the classical localization, is not even defined.
\end{example}


The stack \(\mathcal X^G\) also admits a description via something called the \emph{universal stabilizer} of the action of $G$ on $\mathcal X$, denoted \(\St_{X,G}\) \cite[Section 4.1]{romagny_AlgebraicitySmoothnessFixed2022}. It is the fiber product
\begin{equation}\label{eqn:universal-stabilizer-diagram}
\begin{tikzcd}
  \St_{\mathcal X,G} \arrow[d]\arrow[r]& \mathcal X \arrow[d,"\Delta"] \\
  G\times_S \mathcal X\arrow[r,swap,"\text{act}\times \pr_2"]&\mathcal X\times \mathcal X
\end{tikzcd}.
\end{equation}
Over a point \(x:U\to \mathcal X\) its objects are
\begin{align*}
  \St_{\mathcal X,G}(U) = \{(g,\alpha) ~\mid~ g\in G(U), \alpha:gx\xrightarrow{\sim} x\},
\end{align*}
with multiplication law \((g,\alpha) \cdot (h,\beta) = (gh, \alpha \circ g\beta)\) and identity element \((1,\id_x)\). The leftmost vertical arrow in diagram \eqref{eqn:universal-stabilizer-diagram} is given by \((g,\alpha) \mapsto g\) over \(x:U\to \mathcal X\) and is a morphism of \(\mathcal X\)-group stacks. Its kernel is the inertia stack \(I_{\mathcal X} := \mathcal X\times_{\mathcal X\times \mathcal X}\mathcal X,\) and thus we get an exact sequence
\begin{equation}
  \label{eqn:inertia-stack-sequence}
  1 \to I_{\mathcal X/S} \to \St_{\mathcal X,G} \to G\times_S \mathcal X
\end{equation}
where over \(x:U\to \mathcal X\) the map \(I_{\mathcal X/S}\to \St_{\mathcal X, G}\) is
\begin{align*}
  (\alpha:x\xrightarrow{\sim} x) \mapsto (1, \alpha).
\end{align*}
We call the cokernel of \(I_{\mathcal X/S}(x)\hookrightarrow \St_{\mathcal X, G}(x)\) the \emph{\(G\)-stabilizer} of \(x\). As an aside, the rightmost map in \eqref{eqn:inertia-stack-sequence} admits a section over \(U\) if and only if \(x\) is in the image of \(\eta(U):\mathcal X^G(U)\to \mathcal X(U)\), and \(x:U\to \mathcal X\) in the image of \(\eta\).
\begin{remark}\label{rl-000A}
\end{remark}



\begin{lemma}\label{rl-000B}
\cite[Lemma 4.1.2]{romagny_AlgebraicitySmoothnessFixed2022}
  \label{lem:fixed-stack-parameterizes-splittings}
  $\mathcal X^{G}$ parameterizes sections of \(\St_{\mathcal X,G}\to G\times_S \mathcal X\).
\end{lemma}
\begin{proof}
  Fix a point \(x:U\to \mathcal X\). A splitting of \(\St_{\mathcal X, G}\times_{\mathcal X}U\to G\times_S\mathcal X\times_{\mathcal X}U\) is determined by a collection \(\{\alpha_g\}_{g\in G(U)}\) satisfying the cocycle condition in diagram \eqref{eqn:cocycle-condition-for-fixed-stack}, which is exactly the data of a point in \(\mathcal X^G(U)\) over \(x\).

  In particular, the image of \(\eta(U):\mathcal X^G(U)\to \mathcal X(U)\) is the collection of points \(x\in \mathcal X(U)\) where sequence \eqref{eqn:inertia-stack-sequence}, base changed to \(U\), admits a group-theoretic section over \(U\).
\end{proof}

The fixed stack \(\mathcal X^G\) and the map \(\eta:\mathcal X^G\to \mathcal X\) can be arbitrarily bad, but crucially, it is controllable for algebraic tori acting on algebraic stacks.

\subsection{Fixed stacks for torus actions}\label{rl-000C}
\begin{lemma}\label{rl-000D}
  \label{lem:torus-actions-are-nice}
  \textit{(Torus actions are nice.)}~
  Let \(S\) be a locally noetherian scheme and \(T\) a split torus over \(S\) acting on an algebraic stack \(\mathcal X\) over \(S\).
  \begin{enumerate}[(i)]
    \item If \(\mathcal X\) is Deligne-Mumford, quasi-separated and locally of finite type over \(S\), then \(\eta:\mathcal X^T\to \mathcal X\) is a closed immersion. In particular \(\mathcal X^T\) is algebraic and \(\eta\) is proper and quasi-compact.
    \item If \(\mathcal X\) has affine, finitely presented diagonal, then \(\mathcal X^T\) is algebraic and \(\eta\) is representable by schemes, separated and locally of finite presentation. In particular \(\eta\) is DM-type.
    \item If \(S = \Spec k\) and \(\mathcal X\) is locally of finite type over \(k\) with quasi-compact separated diagonal and affine stabilizers, then \(\mathcal X^T\) is algebraic \cite[Thm. A.20, Rmk A.19]{aranha-khan-latyntsev-etal_VirtualLocalizationRevisited2025} and \(\eta\) is representable by algebraic spaces.
  \end{enumerate}
\end{lemma}
\begin{proof}
$ $
  \begin{enumerate}[(i)]
    \item \cite[Prop A.23]{aranha-khan-latyntsev-etal_VirtualLocalizationRevisited2025}
    \item \cite[Thm 1.2.1 (1)]{romagny_AlgebraicitySmoothnessFixed2022}
          \item The derived stack \(\mathcal X^{hT}\) is 1-Artin by \cite[Thm. A.20]{aranha-khan-latyntsev-etal_VirtualLocalizationRevisited2025}, so its classical truncation \(\mathcal X^T\) is an algebraic stack by Remark \ref{rmk:algebraicity-is-representability} below.
  \end{enumerate}
\end{proof}

\begin{remark}\label{rl-000E}\label{rmk:algebraicity-is-representability}
  The representability of \(\eta\) by algebraic stacks is equivalent to the algebraicity of \(\mathcal X^T\). To see this, first notice that the map \(\eta:\mathcal X^T\to \mathcal X\) is always faithful. Indeed, a morphism \(\phi:(x,\{\alpha_g\})\to (y,\{\beta_g\})\) in \(\mathcal X^T\) is by definition a morphism \(\phi:x\to y\) in \(\mathcal X\) which satisfies the equivariance condition \(\phi\circ \alpha_g = \beta_g \circ (g\cdot \phi)\). This is a condition on \(\phi\), not additional data, so
  \[\Hom_{\mathcal X^T}\big((x,\{\alpha_g\}), (y,\{\beta_g\})\big) \hookrightarrow \Hom_{\mathcal X}(x,y)\]
  is injective. Next, a morphism \(f:\mathcal X\to \mathcal Y\) of algebraic stacks is faithful if and only if it is representable by algebraic spaces \cite[\href{https://stacks.math.columbia.edu/tag/04Y5}{Tag 04Y5}]{stacks-project}, and a stack admitting a morphism representable by algebraic spaces to an algebraic stack is algebraic \cite[\href{https://stacks.math.columbia.edu/tag/03YS}{Tag 03YS}]{stacks-project} hence the algebraicity of \(\mathcal X^T\) is equivalent to the representability of \(\eta\).
\end{remark}

Furthermore, the hypothesis that \(\mathcal X\) have affine diagonal cannot simply be dropped. Example \ref{example:affine-diagonal-is-necessary} is a reconstruction of an example mentioned by Romagny following~\cite[Theorem 1.2.2]{romagny_AlgebraicitySmoothnessFixed2022} found in \cite[Exp. IX, Rem. 7.4]{SGA3-2},~ illustrating that \(\mathcal X^T\) may indeed fail to be algebraic when \(\mathcal X\) does not have affine diagonal. We do not know what minimal hypotheses are necessary, but affine stabilizers are sufficient by \cite[Theorem A.20]{aranha-khan-latyntsev-etal_VirtualLocalizationRevisited2025}.
\begin{lemma}\label{rl-000F}\label{lem:two-term-truncation-of-fixed-locus-cotangent-complex-no-fixed-part}
  If \(T\) is a diagonlizable group over a field which acts on an algebraic stack of finite type over \(k\) such that \(\mathcal X^T\to \mathcal X\) is representable, then the two term truncation of \(L_{\mathcal X^T/\mathcal X}\) has no fixed part with respect to the natural \(T\)-action.
\end{lemma}

\red{This is the proof nearly verbatim from Tom's original notes. NOT COMPLETED (obviously)}
\begin{proof}
  Let \(L = \tau_{\geq -1}L_{\mathcal X^T/\mathcal X}\). First we show the vanishing \(h^0(L^f) = \Omega^f_{\mathcal X^T/\mathcal X}\), for which it suffices to check geometric points. Given \(p\in \mathcal X(k')\), the fiber of \(\mathcal X^T\) over \(p\) is the space of splittings of 
\end{proof}
\red{The following ``proof'' is something I hacked together with Claude after looking at \cite[Corollary A.31]{aranha-khan-latyntsev-etal_VirtualLocalizationRevisited2025}. It's not vetted -- it's an alternate idea which needs verification. Abandoning is also viable.}
\begin{proof}
  This may be proven directly using a Picard-stack argument, but instead we appeal to \cite[Cor. A.31]{aranha-khan-latyntsev-etal_VirtualLocalizationRevisited2025}, translating their derived stack results to classical Artin stacks.
  Let \(\varepsilon:\mathcal X^{hT}\to \mathcal X\) be the derived homotopy fixed stack \cite[Def. A.14, Rem A.15]{aranha-khan-latyntsev-etal_VirtualLocalizationRevisited2025} which is \(T\)-equivariant for the trivial action on the source. For a classical \(S\)-scheme \(U\), one has
  \begin{align*}
    \operatorname{Maps}_S(U,\mathcal X^{hT}) = \operatorname{Maps}_S^T(U, \mathcal X) = \Hom_{T-\St}(\iota(U), \mathcal X) = \mathcal X^T(U),
  \end{align*}
  so the classical truncation of \(\mathcal X^{hT}\) is \(\mathcal X^T\) and \(\eta = \varepsilon \circ \iota\) with \(\iota:\mathcal X^T\to \mathcal X^{hT}\) the classical truncation morphism. By \cite[Rem. A.19]{aranha-khan-latyntsev-etal_VirtualLocalizationRevisited2025} \(BT\) is formally proper over \(S\) and hence we may apply \cite[Cor. A.31]{aranha-khan-latyntsev-etal_VirtualLocalizationRevisited2025} to obtain \(L_{\mathcal X^{hT}/\mathcal X}\simeq (\varepsilon^*L_{\mathcal X})^{mov}[1]\), whose fixed part vanishes. Decomposition of \(D_{qc}(-\times BT)\) into weight-spaces is functorial for base change, so pullback along \(i\) commutes with the fixed/moving decomposition and thus \((i^*L_{\mathcal X^{hT}/\mathcal X})^{f} = 0\). The transitivity triangle for \(\mathcal X^T\xrightarrow{i}\mathcal X^{hT}\xrightarrow{\varepsilon}\mathcal X\) is
  \begin{align*}
    i^*L_{\mathcal X^{hT}/\mathcal X} \to L_\eta \to L_{\mathcal X^T/\mathcal X^{hT}} \xrightarrow{+1},
  \end{align*}
  which is a triangle in \(D_{qc}(\mathcal X^{T}\times_SBT)\). Since \(i\) is the inclusion of the classical truncation, \(L_{\mathcal X^T/\mathcal X^{hT}}\) is 2-connective, meaning \(h^j(\mathcal X^T/\mathcal X^{hT}) = 0\) for \(j \geq -1\). The long exact sequence in cohomology therefore gives \(h^0(i^*L_{\mathcal X^{hT}/\mathcal X}) \simeq h^0(L_\eta)\) and a surjection \(h^{-1}(L_{\mathcal X^{hT}/\mathcal X}) \twoheadrightarrow h^0(L_\eta)\). Taking fixed parts yields the claim.
\end{proof}
We end this section describing how, given an obstruction theory for \(\pi:M\to \mathcal X\), we obtain a canonical obstruction theory for \(M^T\to \mathcal X^T\).
\begin{lemma}\label{rl-000G}\label{lem:obstruction-theory-for-the-fixed-locus}
  Suppose \(E\) is a \(T\)-equivariant relative perfect obstruction theory for \(\pi:M\to \mathcal X\) where \(M\) is a Deligne-Mumford stack and \(\mathcal X\) is an algebraic stack. That is, \(E\) is an element in the \(T\)-equivariant derived category of \(M\) of perfect amplitude \([-1,0]\) together with a \(T\)-equivariant morphism \(\phi: E\to L_{M/\mathcal X}\) which is an obstruction theory in the sense of \cite[Definition 4.4]{behrend-fantechi_IntrinsicNormalCone1997}. If \(\mathcal X^T\) is an algebraic stack, then \((E|_{M^T})^f\) is a canonical relative perfect obstruction theory for \(M^T\) over \(\mathcal X^T\).
\end{lemma}
\begin{proof}
  We have a commutative square (not Cartesian)
  \begin{center}
    \begin{tikzcd}
      M^T \arrow[r, "\iota"] \arrow[d, "\pi^T"] & M \arrow[d, "\pi"] \\
      \mathcal X^T \arrow[r, "\eta"] & \mathcal X
    \end{tikzcd}.
  \end{center}
  Traversing from \(M^T\) to \(\mathcal X\) through \(M\) and through \(\mathcal X^T\) each gives us a transitivity triangle of cotangent complexes, and pulling back \(\phi\) by \(\iota^*\) gives us a morphism \(\iota^*\phi:\iota^*E\to \iota^*L_{M/\mathcal X}\). We can arrange this data as follows:
  \begin{center}
    \begin{tikzcd}
      \iota^*E \arrow[d,"\iota^*\phi "] & & & \\
    \iota^*L_{M/\mathcal X} \arrow[r,"\alpha"] & L_{M^T/\mathcal X} \arrow[r] \arrow[d, equal] & L_{M^T/M} \arrow[r, "+1"] & \dots\\
    (\pi^T)^* L_{\mathcal X^T/\mathcal X} \arrow[r] & L_{M^T/\mathcal X} \arrow[r, "\beta"] & L_{M^T/\mathcal X^T} \arrow[r, "+1"] & \dots
  \end{tikzcd}
  \end{center}
  where the middle equality follows from the commutativity of the diagram. Write \(\iota^*E\) as \(E|_{M^T}\) to match the notation in the statement and define \(\psi\) to be the composition of \(\iota^*\phi\), \(\alpha\) and \(\beta\) after taking fixed parts:
  \begin{align*}
    \psi: \left(E|_{M^T}\right)^f\xrightarrow{\iota^*\phi^f} \left(\iota^*L_{M/\mathcal X}\right)^f \xrightarrow{\alpha^f} \left(L_{M^T/\mathcal X}\right)^f \xrightarrow{\beta^f} \left(L_{M^T/\mathcal X^T}\right)^f.
  \end{align*}
  We show that \(\psi\) is a perfect obstruction theory. First consider the map \(\iota^*\phi\). Right exactness of derived pullback implies that \(h^0(\iota^*E) = \iota^*h^0(E)\) and \(h^0(\iota^*L_{M/\mathcal X}) = \iota^*h^0(L_{M/\mathcal X})\), and \(\iota^*\) preserves isomorphisms, so \(h^0(\iota^*\phi)\) is an isomorphism. For \(h^{-1}(-)\), we get a Tor correction and a diagram with exact rows:
  \begin{center}
  \begin{tikzcd}
    \iota^*h^{-1}(E) \arrow[r] \arrow[d,"\iota^*h^{-1}(\phi)"] & h^{-1}(\iota^*E) \arrow[r]\arrow[d,"h^{-1}(\iota^*\phi)"] & \Tor_1(h^0(E)) \arrow[r] \arrow[d,"\Tor_1(h^0(\phi))"] & 0 \arrow[d] \\
    \iota^*h^{-1}(L_{M/\mathcal X}) \arrow[r] & h^{-1}(\iota^*L_{M/\mathcal X}) \arrow[r] & \Tor_1(h^0(L_{M/\mathcal X})) \arrow[r] & 0 \\
  \end{tikzcd}
  \end{center}
  \(\Tor\) preserves isomorphisms so \(\Tor_1(h^0(\phi))\) is an isomorphism. The morphism \(h^{-1}(\phi)\) is surjective by assumption and \(\iota^*(-)\) is right exact, so \(\iota^*h^{-1}(\phi)\) is also surjective. Applying \cite[Lemma 05QA]{stacks-project} then gives us that \(h^{-1}\iota^*\phi\) is surjective as well.

  Now consider the map \(\alpha\). The fixed portion of the long exact sequence of the transitivity triangle involving \(\alpha\) reads
  \begin{align*}
    &h^{-1}(\iota^*L_{M/\mathcal X})^f \xrightarrow{h^{-1}(\alpha^f)}h^{-1}(L_{M^T/\mathcal X})^f \to h^{-1}(L_{M^T/M})^f \to \dots\\\dots \to ~& h^{0}(\iota^*L_{M/\mathcal X})^f \xrightarrow{h^{0}(\alpha^f)}h^{0}(L_{M^T/\mathcal X})^f \to h^{0}(L_{M^T/M})^f \to 0
  \end{align*}
  

  Since \(\iota\) is a closed immersion (Lemma \ref{lem:torus-actions-are-nice}) \(h^0(L_{M^T/M}) = \Omega_{M^T/M} = 0\), and applying Lemma \ref{lem:two-term-truncation-of-fixed-locus-cotangent-complex-no-fixed-part} to \(\mathcal M\) gives us \(h^{-1}(L_{M^T/M})^f = 0\). This implies \(h^0(\alpha^f)\) is an isormophism and \(h^{-1}(\alpha^f)\) is surjective.

  Finally consider \(\beta\). Because \(T\) acts trivially on \(\mathcal M^T\) and \(\mathcal X^T\), \(L_{\mathcal M^T/\mathcal X^T}\) is of pure weight zero and the weight-zero part of the long exact sequence of the bottom triangle reads
  \begin{align*}
    &h^{-1}(\iota^*L_{\mathcal X^T/\mathcal X})^f \to h^{-1}(L_{M^T/\mathcal X})^f \xrightarrow{h^{-1}(\beta^f)} h^{-1}(L_{M^T/\mathcal X^T})^f \to \dots\\\dots \to ~& h^{0}(\iota^*L_{\mathcal X^T/\mathcal X})^f \to h^{0}(L_{M^T/\mathcal X})^f \xrightarrow{h^{0}(\beta^f)} h^{0}(L_{M^T/\mathcal X^T})^f \to 0.
  \end{align*}
  Since \((\pi^T)^*\) is right \(t\)-exact and compatible with weights, \(h^0((\pi^T)^*L_{\mathcal X^T/\mathcal X})^f = (\pi^T)^*(\Omega^f_{\mathcal X^T/\mathcal X}) = 0\) by Lemma \ref{lem:two-term-truncation-of-fixed-locus-cotangent-complex-no-fixed-part}. Hence \(h^0(\beta^f)\) is an isomorphism and \(h^{-1}(\beta^f)\) is surjective. Thus \(h^0(\psi)\) is the composite of isomorphisms and \(h^{-1}(\psi)\) is the composite of surjections, so we conclude that \(\psi\) is a perfect obstruction theory.
\end{proof}


\subsection{Non-algebraicity of the fixed stack}\label{rl-000H}\label{sec:non-algebraicity-of-the-fixed-stack}

\begin{example}\label{rl-000I}\label{example:affine-diagonal-is-necessary}
  Let \(R = \mathbb C[\hspace{-2pt}[q]\hspace{-2pt}]\), $K = \mathbb C(\hspace{-2pt}(q)\hspace{-2pt})$ \(S = \Spec R\) and let \(W \to S\) be the Weierstrass model
  \begin{align*}
    y^2 + xy = x^3 + a_4x + a_6
  \end{align*}
  of the Tate curve \cite{silverman_AdvancedTopicsArithmetic1999}. Now let \(E\) be the smooth locus of \(W\). Then \(E\) is a quasi-compact open subscheme of the projective \(S\)-scheme \(W\), hence \(E\to S\) is separated, finitely presented and smooth. Its generic fiber is an elliptic curve and its central fiber is the smooth locus of a nodal cubic, so \(E_{\mathbb C} = \mathbb G_m\).

  Set \(\mathcal X = BE = [S/E]\) and let \(\mathbb G_m\) act trivially on \(\mathcal X\). Then of the hypotheses in \cite[Theorem 1.2.1]{romagny_AlgebraicitySmoothnessFixed2022}, \(\mathcal X\) only fails the affine diagonal requirement. Indeed, \(\mathcal X\) is an algebraic stack by \cite[06PL]{stacks-project} and the group \(\mathbb G_m\) is reductive. However, if we base change \(\Delta_{BE}:BE\to BE_SBE\) along \(S\to BE\times_SBE\) given in each component by the trivial torsor \(S\to BE\), then we get
  \begin{align*}
    BE\times_{BE\times_SBE} S \cong E.
  \end{align*}
  The fiber of $E$ over the generic point of \(S\) is the proper curve \(E_K\to \Spec K\), so \(E_K\) is not affine and neither is \(\Delta_{BE}\).
  
  We claim that \((BE)^{\mathbb G_m}\) is not algebraic. The Romagny fixed stack is
  \begin{align*}
    X^{\mathbb G_m} \cong BE\times_S F, \hspace{1em} F = \underline{Hom}_{S-gp}(\mathbb G_m, E),
  \end{align*}
  and \(\eta:X^{\mathbb G_m}\to X\) is the projection to the first factor by the description at the end of Remark \ref{rmk:description-of-homotopy-fixed-stack}. Since \(\eta\) is faithful (Remark \ref{rmk:algebraicity-is-representability}) it is representable by algebraic spaces by \cite[\href{https://stacks.math.columbia.edu/tag/04Y5}{Tag 04Y5}]{stacks-project}, so if \(\mathcal X^{\mathbb G_m}\) were algebraic then its base change \(F = S\times_{BE}\mathcal X^{\mathbb G_m}\) along the trivial torsor would be an algebraic space over \(S\). We show that no algebraic space can have the values \(F\) has by exhibiting a compatible family \(f_n \in F(A_n)\) for \(A_n \in R/(q^n)\) with \(f_1 \neq 0\), while \(F(R) = 0\).
  
  This is straightforward. Observe first that any group homomorphism \(f:\mathbb G_{m,R}\to E\) must induce a group homomorphism \(f_K:\mathbb G_{m,K}\to E_K\). This extends to a morphism \(\mathbb P^1_K\to E_K\), but the only such morphism is constant, hence \(f_K\) is trivial and \(F(R) = 0\). If we instead reduce to \(\mathbb C = R/(q)\), then we get
  \[F(\mathbb C) = \Hom(\mathbb G_m, \mathbb G_m) = \mathbb Z.\]
  Next we need to show there exists some non-trivial element in \(\lim \Hom(\mathbb G_{m,A_n}, E_{A_n})\) for \(A_n = R/(q^n)\). One can do this directly by showing the identity morphism \(f_1:\mathbb G_{m,\mathbb C} \to \mathbb G_{m,\mathbb C}=E_{\mathbb C}\) lifts to a morphism \(f_n:\mathbb G_{m,A_n}\to E_{A_n}\) where \(A_n = R/(q^n)\). In homogeneous coordinates, it's given by
  \begin{align*}
    f_n(t) = \big[t(1-t) + (1 - t)^3P_n(t,q) ~ : ~ t^2 + (1 - t)^3 Q_n(t,q) ~ : ~ (1 - t)^3\big]
  \end{align*}
  for certain Laurent polynomials \(P_n,Q_n \in qA_n[t,t^{-1}]\) obtained by truncating the Tate series modulo \(q^{n}\).

  Suppose now that $\mathcal{X}^{\mathbb{G}_m}$ is algebraic, so that $F$ is an algebraic space over $S$, and choose an \'etale surjection $p:V\to F$ from a scheme. The $\mathbb{C}$-point $f_1\in F(A_1)$ lifts to a $\mathbb{C}$-point $v$ of $V$, since $p$ is surjective and $\mathbb C$ is algebraically closed. As $p$ is \'etale, hence formally \'etale, and $\Spec A_{n-1}\subset \Spec A_n$ is a nilpotent thickening, the maps $f_n:\Spec A_n\to F$ lift uniquely and compatibly to $\lambda_n:\Spec A_n\to V$ with $\lambda_1=v$. A local Artinian scheme lands in every open neighbourhood of its closed point, so all $\lambda_n$ factor through a single affine open $\Spec B\ni v$ of $V$; they are then given by compatible ring maps $B\to A_n$, that is, by a ring map $B\to\varprojlim A_n=R$. The resulting composite $\Spec R\to V\to F$ is an element $f\in F(R)$ with $f|_{A_n}=f_n$ for every $n$. But $F(R)=0$, so $f_1=0$, contradicting $f_1=\mathrm{id}_{\mathbb{G}_m}$. Hence $\mathcal{X}^{\mathbb{G}_m}=(BE)^{\mathbb{G}_m}$ is not algebraic.
\end{example}


\begin{remark}\label{rl-000J}\label{rmk:effectivity-for-algebraic-spaces}
  The arguement above uses only that \'etale morphisms have a unique infinitesimal lifting, that schemes satisfy effectivity and that \(R\) is \(q\)-adically complete. It therefore shows that condition (F3) of \cite[Theorem 2.3.1]{romagny_AlgebraicitySmoothnessFixed2022} holds for every \emph{algebraic space}, which is why its failure obstructs algebraicity and not merely representability by an unramified spearated \emph{scheme}.
\end{remark}


\subsection{Remaining facts from Tom's note that I haven't fit in anywhere else yet so they get their own subsection}\label{rl-000K}


\begin{proposition}\label{rl-000L}\label{prop:condition-for-monomorphism}
  If $G$ is any connected group scheme over a field acting on a stack $X$ and \(I_X\to X\) is unramified, the map $X^{G}\to X$ is a monomorphism.
\end{proposition}
\begin{proof}
  Suppose we have a map $h:A\to X\in X(A)$ with two lifts $(f,g:A\to X^{G})$ under the natural map $\epsilon:X^{G}\to X$, i.e. that there is an isomorphism $\alpha:f\Rightarrow g$. Since $X^{G}$ parameterizes splittings of the map $\St_{X,G}\to G_X$ (Lemma~\ref{lem:fixed-stack-parameterizes-splittings}) this is the same as having two sections $s_f$ and $s_g$ of $\alpha_A:h^*\St_{X,G}\to G_A$. Since both \(s_f\) and \(s_g\) are sections, composing with \(\alpha_A\) defines the same map. Therefore they differ by some element of \(h^*I_X\) by exactness in~\eqref{eqn:inertia-stack-sequence}; that is, the image of $s_fs^{-1}_g$ lands in the image of the inclusion $h^*I_X\hookrightarrow h^*\St_{X,G}$ where $s_g^{-1}$ is the map $G_A\to h^*\St_{X,G}$ defined by $x\mapsto (s_g(x))^{-1}$. The image of \(s_fs^{-1}_g\) must be connected because \(G\) is connected, and it must land in the identity component of \(h^*I_X\) since this is a homomorphism of group algebraic spaces. However the identity component consists only of the identity itself since \(I_X \to X\) is unramified, hence \(s_fs^{-1}_g\) is trivial.
\end{proof}


The next proposition Graber states provides a properness criterion for the morphism \(\epsilon:X^{G}\to X\).

\begin{proposition}\label{rl-000M}
  If \(G\) is any smooth connected group scheme over a field \(k\), and \(G\) acts on a quasi-separated algebraic stack \(X\) whose diagonal is finite type and unramified, then the morphism \(X^{G}\to X\) satisfies the valuative criterion for properness.
\end{proposition}


An \(S\)-stack \(X\) is Deligne-Mumford precisely when it has an unramified diagonal morphism \(\Delta:X\to X\times_S X\). This means we can trade out this proposition for the following statement proven in  \cite{aranha-khan-latyntsev-etal_VirtualLocalizationRevisited2025}, which requires only that \(X\) be quasi-DM, or in other words, that \(\Delta\) is locally quasi-finite.

\begin{proposition}\label{rl-000N}
  Suppose $S$ is a locally Noetherian scheme and \(G\) is a group scheme with smooth and connected fibers over \(S\). Let \(X\) be a quasi-DM stack with a \(G\)-action which is locally of finite type over \(S\). Then the canonical morphism \(\epsilon:X^{G}\to X\) is essentially proper (i.e. \ it locally of finite type and satisfies the valuative criterion for properness.)
\end{proposition}
\begin{proof}
  \cite[Proposition A.22]{aranha-khan-latyntsev-etal_VirtualLocalizationRevisited2025}
\end{proof}


\red{This lemma and the next should probably be in the ``reconciliation of different fixed stacks'' section, if it exists}
\begin{lemma}\label{rl-000O}
  Suppose $p\in X^{T}(k')$ is a geometric point, then the morphism $\epsilon:X^{T}\to X$ is an isomorphism onto $X^T_{naive}$ in a formal neighborhood of $p$.
\end{lemma}
\begin{proof}
  We know $\epsilon$ is a monomorphism by Proposition~\ref{prop:condition-for-monomorphism}. \isaac{Graber says all we need here is the fact that, given a $T$-equivariant morphism from a scheme $f:Z\to X$ and a thickening $\overline Z$ of $Z$, every $T$-invariant extension of $f$ to $\overline Z$ lifts to a $T$-equivariant extension.}
\end{proof}

\begin{lemma}\label{rl-000P}\label{lem:finite-cover}
  If $X$ is a DM stack of finite type over a field with an action of a torus $T$, then there exists a finite covering of tori $T'\to T$ such that the natural morphism $X^{T'} \to X^{T'}_{naive}$ is surjective on geometric points.
\end{lemma}


\input{nodes/rl-000Q}
\input{nodes/rl-0015}
\input{nodes/rl-001A}
\end{document}
